%% A small CCC recovered from its named full equational diagram.

\subsection{Recovering a category from its internal language}
\label{sec:p10-internal-cartesian-language}

Fix a small cartesian closed category \(\mathcal C\) with given terminal
object \(1\), binary products, and exponentials \(B^A\) with evaluation
\(\operatorname{ev}_{A,B}:B^A\times A\to B\).
Small means that its objects and arrows form sets in the chosen universe.
The chosen universal constructions are part of the data.
We construct a particular equational language for this category.

\begin{hbdefinition}{def:p10-internal-ccc-language}{The named internal language of a small CCC}
For each object \(A\) introduce a base type \(\underline A\), using
distinct formal symbols for distinct objects.  Form all other types by
the arrow constructor of Definition~\ref{def:p10-stlc-beta-eta}.
Use the same variables, application, and typed abstraction, and add
the following closed constants:
\[
\begin{array}{ll}
 c_f:\underline A\mathbin{\to_{\mathrm{ty}}}\underline B
       &(f:A\to B),\\
 u:\underline1,&\\
 \mathsf{pair}_{A,B}:\underline A\mathbin{\to_{\mathrm{ty}}}
        \underline B\mathbin{\to_{\mathrm{ty}}}\underline{A\times B},&\\
 e_{A,B}:\underline{B^A}\mathbin{\to_{\mathrm{ty}}}
        (\underline A\mathbin{\to_{\mathrm{ty}}}\underline B),&\\
 l_{A,B}:(\underline A\mathbin{\to_{\mathrm{ty}}}\underline B)
        \mathbin{\to_{\mathrm{ty}}}\underline{B^A}.&
\end{array}
\]
Each constant has its declared type in any context; substitution leaves
it fixed.  No product or unit type constructor is added.
The structural constants compare named objects with contexts and arrow
types; \(\underline{B^A}\) and
\(\underline A\mathbin{\to_{\mathrm{ty}}}\underline B\) are distinct
type expressions.

Interpret \(\underline A\) by \(A\), and write \(A_T=\llbracket T\rrbracket\)
for the recursively interpreted object of a type \(T\).
As in Definition~\ref{def:p10-stlc-ccc-interpretation}, put
\[
 P_{()}=1,\qquad P_{\Gamma\cdot(T)}=P_\Gamma\times A_T.
\]
In particular, \(P_{(\underline A)}=1\times A\), not \(A\).
Terms and substitutions use the categorical interpretation of that
definition, extended by the declared constants.

To specify a closed constant of function type \(S\mathbin{\to_{\mathrm{ty}}}T\)
by an arrow \(a:A_S\to A_T\), use the closed denotation
\(\Lambda(a\circ\pi_2):1\to A_T^{A_S}\), where
\(\pi_2:1\times A_S\to A_S\).
In this sense \(c_f\) denotes \(f\), while \(e_{A,B}\) and
\(l_{A,B}\) denote the identity of \(B^A\).
The constant \(u\) denotes \(1_1:1\to1\).
The denotation of \(\mathsf{pair}_{A,B}\) is obtained by first
transposing in \(B\), then in \(A\), the arrow
\[
 ((1\times A)\times B)\longrightarrow A\times B,
 \qquad \langle\pi_2\circ\pi_1,\pi_2\rangle.
\]
Thus its applications denote categorical pairing.
The signature is a set because \(\mathcal C\) is small.
\index{internal language!cartesian closed category}
\end{hbdefinition}

\begin{hbdefinition}{def:p10-internal-ccc-diagram}{The full true equational diagram}
Let \(\operatorname{Diag}(\mathcal C)\) contain every well-typed equation
\[
 \Theta_\Gamma\vdash M=N:T
\]
whose two denotations are equal arrows \(P_\Gamma\to A_T\).
This is a set of equations-in-context.
Let \(=_{\mathcal C}\) be the least typed equational theory containing
this diagram and the \(\beta\eta\) equations, closed under equivalence,
application and abstraction congruence, weakening, renaming, and typed
capture-avoiding simultaneous substitution.
Define \(\mathcal S_{\mathcal C}\) by the context and tuple construction
of Definition~\ref{def:p10-stlc-context-category}, now using these
constants and componentwise \(=_{\mathcal C}\) equality.

Some equations in the diagram are listed below. Each line is read in
its own typing context: \(f,g\) are composable, \(!_A:A\to1\), the
product projections have domain \(A\times B\), and their variables
have the corresponding named types.
\[
\begin{gathered}
 c_{1_A}x=x,\qquad c_g(c_fx)=c_{g\circ f}x,\qquad c_{!_A}x=u,\\
 c_{\pi_1}(\mathsf{pair}_{A,B}xy)=x,\qquad
 c_{\pi_2}(\mathsf{pair}_{A,B}xy)=y,\\
 \mathsf{pair}_{A,B}(c_{\pi_1}z)(c_{\pi_2}z)=z
                     \quad(z:\underline{A\times B}),\\
 e_{A,B}(l_{A,B}h)=h,\qquad l_{A,B}(e_{A,B}z)=z,\\
 (e_{A,B}z)a
   =c_{\operatorname{ev}_{A,B}}(\mathsf{pair}_{B^A,A}za).
\end{gathered}
\]
In the exponential lines,
\(h:\underline A\mathbin{\to_{\mathrm{ty}}}\underline B\),
\(z:\underline{B^A}\), and \(a:\underline A\).
Terminal uniqueness is \(z=u\) in context \(z:\underline1\).
For \(f:D\times A\to B\), the diagram also records transposition:
\[
 (e_{A,B}(c_{\Lambda f}d))a
       =c_f(\mathsf{pair}_{D,A}da).
\]
These displayed equations describe part of the diagram; the theory
contains all true equations-in-context, not just this subfamily.
No effective enumeration or finite axiomatization of that theory is
asserted.
\index{equational diagram!internal language}
\end{hbdefinition}

The naming and structural equations follow the internal-language method
of \cite[\S2.3, pp.~10--11, Proposition~2.3.1]{Pitts1995Draft}.
The full semantic diagram is the construction in \S4.3, pp.~35--36;
the context category and its type structure are in \S\S4.2--4.4,
pp.~30--38.  These locators refer to the draft dated 18 May 1995.

\begin{hbtheorem}{thm:p10-internal-ccc-comparisons}{proved=annex:p10.internal.1}{Named type and context comparisons}
For every type \(T\), there are closed terms
\[
 \kappa_T:T\mathbin{\to_{\mathrm{ty}}}\underline{A_T},
 \qquad
 \iota_T:\underline{A_T}\mathbin{\to_{\mathrm{ty}}}T
\]
such that applying either comparison to a well-typed argument preserves
that argument's denotation in every context.
Their two composites equal identity terms in \(=_{\mathcal C}\).
For every context \(\Gamma\), there are inverse context arrows
\[
 \operatorname{pack}_\Gamma:\Gamma\longrightarrow(\underline{P_\Gamma}),
 \qquad
 \operatorname{unpack}_\Gamma:(\underline{P_\Gamma})\longrightarrow\Gamma.
\]
The component packing term denotes \(1_{P_\Gamma}\), while the packing
\emph{context arrow} denotes
\(\langle!,1\rangle:P_\Gamma\to1\times P_\Gamma\).
The unpacking context arrow denotes
\(\pi_2:1\times P_\Gamma\to P_\Gamma\).
These comparisons include the empty context.
\end{hbtheorem}

\begin{hbtheorem}{thm:p10-internal-ccc-equivalence}{proved=annex:p10.internal.2}{Evaluation recovers the given category}
Evaluation defines a full and faithful functor
\[
 E:\mathcal S_{\mathcal C}\longrightarrow\mathcal C,
 \qquad E(\Gamma)=P_\Gamma.
\]
It is essentially surjective.
An explicit quasi-inverse \(J\) sends \(A\) to the singleton context
\((\underline A)\) and \(f:A\to B\) to the term tuple \((c_fx)\).
The unit is \(\eta_\Gamma=\operatorname{pack}_\Gamma\), and the counit is
\[
 \varepsilon_A=\pi_2:EJ(A)=1\times A\longrightarrow A.
\]
They are natural isomorphisms and satisfy both triangle identities.
\end{hbtheorem}

\begin{hbtheorem}{thm:p10-internal-ccc-structure}{proved=annex:p10.internal.3}{Preservation of cartesian closed structure}
The category \(\mathcal S_{\mathcal C}\) is cartesian closed, using
the empty context, concatenation, and the curried-list exponential
\(\Delta^\Gamma\) of Theorem~\ref{thm:p10-stlc-context-ccc}.
Evaluation preserves the terminal object and has the canonical product
comparison isomorphism
\[
 j_{\Gamma,\Delta}:P_{\Gamma\cdot\Delta}\longrightarrow
 P_\Gamma\times P_\Delta,
 \qquad
 j_{\Gamma,\Delta}=\langle E(\pi_1),E(\pi_2)\rangle.
\]
If \(\Omega=\Delta^\Gamma\) and
\(\operatorname{ev}_{\rm syn}:\Omega\cdot\Gamma\to\Delta\) is its
syntactic evaluation, the exponential comparison is
\[
 \chi_{\Gamma,\Delta}
  =\Lambda\bigl(E(\operatorname{ev}_{\rm syn})
                   \circ j_{\Omega,\Gamma}^{-1}\bigr):
 P_\Omega\longrightarrow P_\Delta^{P_\Gamma}.
\]
It is an isomorphism.  Thus the equivalence preserves cartesian closed
structure up to these specified comparisons, including empty contexts.
\end{hbtheorem}

This reconstructs the given category from its named true equational
theory.  The bare \(\beta\eta\) context category instead has the
universal semantic property of
Theorem~\ref{thm:p10-stlc-ccc-completeness}; it need not be equivalent
to an arbitrary category in which its terms can be interpreted.
